\minitoc

Maintenant que nous avons traité les principaux problèmes sur lesquels 
on peut tomber en considérant des connaissances, attardons-nous sur l'inconsistance. 
En effet, en logique porpositionnelle, dès lors qu'une base contient une contradiction, 
tout et son contraire peuvent être déduit. 
Or, en pratique, il se peut que nous devions traiter une base inconsistante. En effet, deux 
capteurs peuvent fournir des données conflictuelles, des experts peuvent avoir des avis divergeants, etc... 
Dans ce cas, il est nécessaire d'adopter des stratégies pour êre capable de raisonner 
de façon "juste" sur ces connaissances. 

Ce chapitre est la réelle continuité du précédent, toutes les notations et définitions seront donc 
réutilisées. 

% ========================================================================================================
% La syntaxe pour restaurer la consistance

\section{Restauration de la consistance avec la syntaxe}

L'idée principale lors de la rencontre d'une base inconsistante $BC$ va être de la séparer 
en deux types de connaissances : celles dont on est sûrs et les autres. 

\begin{definition}[Base à Noyau]
    On appellera \emph{base de connaissances à noyau} une base $BC$ telle que $ BC = (N, H)$ où :
    \begin{itemize}
        \item $N$ est le \textbf{noyau} de la base. Ce sont l'ensemble des connaissances 
        sûres de la base. Il est nécessairement \emph{consistant}. 
        \item $H$ sont les \textbf{croyances} de la base. Les connaissances qu'il contient 
        peuvent être sujettes à révision. 
    \end{itemize} 
\end{definition}

\begin{example}
    La base $BC = \{ \text{Oiseau}, \lnot \text{Vole}, \text{Vole}\}$ est inconsistante. 
    On peut donc la scinder en deux parties : son noyau et les croyantes. Ainsi : 
        $$ N = \{ \text{Oiseau} \} \quad \text{et} \quad H = \{ \lnot \text{Vole}, Vole \} $$
\end{example}

Or les croyances d'une base ne doivent pas être ignorées pour autant. 
Il convient alors de définir une extension possible des noyaux de base avec une partie
des croyances. 

\begin{definition}[Scenario]
    Soit $BC = (N, H)$ une base à noyau. On appelle $S \subseteq BC$ un \emph{scénario} de $BC$ si : 
    \begin{enumerate}[label=\roman*)]
        \item $ N \subseteq S$
        \item $S$ est consistant. 
        \item $S$ est maximal pour l'inclusion ensembliste. 
            $$ i.e \not \exists S' \subseteq BC \text{ tel que } S' \text{ est consistant et } S \subset S' $$ 
    \end{enumerate}
\end{definition}

De tels scénarios vont donc nous permettre d'inférer différentes nouvelles connaissances de 
la base $BC$, caractérisons-en quelques-unes. 

\begin{definition}[Conséquences]
    Soit $BC = (N, H)$. On définit différents types de connaissances $ \varphi$. 
    \begin{itemize}
        \item $ \varphi$ est une connaissance \emph{universelle (ou forte)} de $BC$ si $ \varphi$ 
        se déduit de tout scénario de $BC$. 
        \item $ \varphi$ est une connaissance \emph{existentielle (ou faible)} de $BC$ si $ \varphi$ 
        se déduit d'au moins un scénario de $BC$. 
        \item $ \varphi$ est une connaissance \emph{argumentative} de $BC$ si $ \varphi$ 
        est une conséquence faible et $ \lnot \varphi$ n'est pas une conséquence faible de $ BC$. 
    \end{itemize}
\end{definition}


% ========================================================================================================
% L'Approche Sémantique : Modèles de Herbrand Préférés

\section{Caractérisation Sémantique}

Plutôt que de manipuler des sous-ensembles de formules, l'approche sémantique compare les \emph{interprétations} 
(mondes possibles) du noyau $N$ en fonction de leur capacité à satisfaire un maximum de croyances de $H$.

\subsection{Relation de préférence entre modèles}

On utilise la notation $\omega \models H$ pour désigner l'ensemble des formules de $H$ qui sont vraies dans 
l'interprétation $\omega$ : $\{ h \in H \mid \omega \models h \}$.

\begin{definition}[Préférence Sémantique]
    Soient $\omega$ et $\omega'$ deux modèles du noyau $N$. On dit que $\omega$ est \emph{préféré} à $\omega'$, 
    noté $\omega \succeq_H \omega'$, si $\omega$ satisfait plus de croyances de $H$ que $\omega'$ au sens de l'inclusion :
        $$ \{ h \in H \mid \omega' \models h \} \subseteq \{ h \in H \mid \omega \models h \} $$
\end{definition}

\begin{definition}[Modèle Préféré]
    Un modèle $\omega$ de $N$ est dit \emph{préféré} s'il est maximal pour la relation $\succeq_H$. Autrement dit, il 
    n'existe aucun autre modèle $\omega''$ de $N$ tel que $\omega'' \succeq_H \omega$.
\end{definition}

\subsection{Lien entre Syntaxe et Sémantique}

Il existe une dualité parfaite entre les scénarios (approches syntaxiques) et les modèles préférés :

\begin{theorem}[Équivalence Scénario/Modèle Préféré]
    Soit $BC = (N, H)$ une base à noyau.
    \begin{itemize}
        \item Tout modèle d'un scénario $S$ est un modèle préféré de $BC$.
        \item Inversement, tout modèle préféré de $BC$ est modèle d'un (et d'un seul) scénario $S$.
    \end{itemize}
\end{theorem}

\begin{remark}
    Une formule $\varphi$ est une conséquence \emph{universelle} de $(N, H)$ si et seulement si elle est satisfaite 
    par tous les modèles préférés de la base.
\end{remark}





